\section{Coordinate Conversion and Vizard Implementation}
\label{sec:conversion}

\subsection{Kinematic Axis and Quaternion Mapping}
The Xsens MVN reference coordinate system establishes $X$ pointing forward, $Y$ lateral to the left, and $Z$ vertical upward. Conversely, the virtual reality graphics engine \Vizard{} (v8.3.1) \cite{worldviz2024vizard} employs a standard right-handed convention with $X$ pointing lateral to the right, $Y$ vertical upward, and $Z$ forward. Consequently, the linear spatial translation and quaternion transformation are uniquely governed by:
\begin{align}
  \mathbf{p}_{\text{Vizard}} & = \begin{bmatrix} -y_{\text{Xsens}} \\ z_{\text{Xsens}} \\ x_{\text{Xsens}} \end{bmatrix}, \label{eq:pos_map}      \\
  \mathbf{q}_{\text{Vizard}} & = \left(-y_{\text{Xsens}},\, z_{\text{Xsens}},\, x_{\text{Xsens}},\, -w_{\text{Xsens}}\right). \label{eq:quat_map}
\end{align}

To eliminate arbitrary horizontal laboratory offsets while strictly preserving true vertical floor height, an origin shift is applied such that the initial pelvis position ($t = 0\text{ s}$) resides precisely at $(0,\, y_{\text{pelvis},0},\, 0)\text{ m}$.

\subsection{Dual Representation: Kinematic Skeleton and Visual Avatar}
To prevent visual mesh artifacts from confounding quantitative biomechanical metrics, the virtual environment incorporates a dual representation:
\begin{enumerate}[leftmargin=*]
  \item \textbf{Segment Skeleton:} Composed of rigid spherical joint markers (\texttt{joint nodes}) and dynamic line segments (\texttt{bones\_lines}) rendered directly from true metric Xsens coordinate positions. Euclidean distances, the racket proxy location, and affordance criterion tests are computed exclusively within this veridical metric space.
  \item \textbf{Retargeted Avatar (\texttt{vcc\_male.cfg}):} Positioned adjacent to the skeleton (lateral offset of $+1.2\text{ m}$ in $X$) and driven by relative sensor rotations applied to the standard bip-skeleton hierarchy (\texttt{Bip01}). This avatar serves qualitative visual inspection of actor volume.
\end{enumerate}

\subsection{Virtual Environment and Scale References}
The virtual scene faithfully reproduces official ITTF tournament table tennis dimensions: length of $2.74\text{ m}$, width of $1.525\text{ m}$, tabletop surface height of $0.76\text{ m}$, and a center net height of $0.1525\text{ m}$. The table automatically aligns along the horizontal swing vector at peak velocity ($t = 3.083\text{ s}$) and sits at a standard clearance distance of $0.40\text{ m}$ from the player's initial pelvis. A 3D physical scale ruler of $1.0\text{ m}$ with alternating black-and-white $10\text{ cm}$ segments is placed on the floor, providing an independent, human-readable optical reference.

\subsection{Interactive Playback Controls}
The scene provides comprehensive interactive keyboard control:
\begin{itemize}[leftmargin=*]
  \item \texttt{SPACE}: Toggle Play / Pause.
  \item \texttt{R}: Restart playback to initial frame ($t = 0\text{ s}$, frame 0).
  \item \texttt{S}: Toggle playback speed between real-time ($1.0\times$) and calibrated slow-motion ($0.25\times$), explicitly labeled in red in the HUD.
  \item \texttt{1}, \texttt{2}, \texttt{3}, \texttt{4}: Select environmental affordance condition (ball location).
  \item \texttt{V}: Show / hide the objective criterion result (blind judgment mode).
  \item \texttt{C}: Cycle through 4 camera vantage points (front-side, front, top-down, and behind player).
  \item \texttt{G}: Step directly to the three validation timestamps.
  \item \texttt{P}: Print instantaneous values (time, frame, flexion, height, ball-racket distance) to console.
  \item \texttt{T}, \texttt{A}, \texttt{K}: Toggle visibility of racket sweep trail, avatar, and skeleton.
\end{itemize}
